\PassOptionsToPackage{dvipsnames}{xcolor}
\documentclass[aspectratio=169,12pt]{beamer}

\usepackage[utf8]{inputenc}
\usepackage[T1]{fontenc}

\usetheme{Rennes}

\usepackage{fontawesome5}
\usepackage{listings}
\usepackage{tikz}

\lstset{
  basicstyle=\ttfamily\small,
  commentstyle=\color{steelGray},
  keywordstyle=\color{lavenderPurple},
  frame=single,
  columns=fullflexible,
  aboveskip=0.5ex,
  belowskip=0.5ex,
}

\lstdefinelanguage{diff}{
  morecomment=[f][\color{lapisBlue}]{@@}, % group identifier
  morecomment=[f][\color{crimsonRed}]-,   % deleted lines
  morecomment=[f][\color{appleGreen}]+,   % added lines
  morecomment=[f][\color{wisteriaViolet}]{---}, % header lines (must appear after +,-)
  morecomment=[f][\color{wisteriaViolet}]{+++},
}

\newcommand{\mono}[1]{\texttt{#1}}
\newcommand{\strong}[1]{\textbf{#1}}

\pgfdeclareimage[height=1.2em]{smalliconarticle}{beamericonarticle}
\newcommand{\bibicon}{\raisebox{-4pt}{\pgfuseimage{smalliconarticle}}}
\newcommand{\bibarticle}[1]{\bibicon\xspace#1\xspace}
\newcommand{\bibfootline}[1]{\gdef\myfootline{\bibarticle{#1}}}


%================================================================================

\title{\mono{Rennes}}
\subtitle{An unofficial Beamer theme}
\author{\strong{Aymeric Blot}}
\institute{Université de Rennes}
\date{2024}

\titlegraphic{%
  \includegraphics[height=5ex]{logo-univ-rennes}\qquad%
  \includegraphics[height=5ex]{logo-irisa}%
}

\begin{document}

%================================================================================
\section{Introduction}
%================================================================================

% define background
\setbeamertemplate{background canvas}{\includegraphics{bg_title}}
% add title slide
\frame[noframenumbering]{
  \titlepage
  \setfootline{\faLink~\url{https://github.com/bloa/beamertheme-rennes}}
}
% remove background for following slides
\setbeamertemplate{background canvas}{}


%================================================================================
\section{Contents}
%================================================================================

\begin{frame}[fragile]{Background}
  \structure{No default background}\\
  The \mono{Rennes} theme does not include any predefined background.\\
  The title page background in this slide deck, showcasing the University of Rennes logo, is only defined locally for this user guide.

  \medskip
  \infoed{\structure{Example}}
  \begin{lstlisting}[
      basicstyle=\ttfamily\footnotesize,
      language={[LaTeX]TeX},
    ]
% define custom background
\setbeamertemplate{background canvas}{\includegraphics{bg_title}}
% add title slide
\frame[noframenumbering]{\titlepage}
% remove background for following slides
\setbeamertemplate{background canvas}{}
  \end{lstlisting}
\end{frame}

\begin{frame}{Accent Colours Commands \warn{(new)}}

  \structure{More than just \mono{\textbackslash alert}}\\
  The \mono{Rennes} theme include \alert{six} different easy-access colour commands:

  \smallskip
  \begin{description}
  \item[\mono{\textbackslash accent}] defaults to the \accent{current \mono{structure} colour}
  \item[\mono{\textbackslash alert}] sharp contrast for \alert{critical information}
  \item[\mono{\textbackslash warn}] highlight for other \warn{important stuff}
  \item[\mono{\textbackslash info}] neutral third colour \info{just for the sake of it}
  \item[\mono{\textbackslash subtle}] very light touch of colour for \subtle{minor stuff}
  \item[\mono{\textbackslash shade}] grey for \shade{disabled stuff}, for example
  \end{description}

  \medskip
  \structure{Difference from Beamer's \mono{\textbackslash alert}}\\
  These six commands only change the \emph{default} font colour.
\end{frame}

\begin{frame}[fragile]{Accent\subtle{ed} Colours Commands \warn{(new)}}
  \structure{\mono{-ed} commands variants:
  \mono{\textbackslash alerted} \mono{\textbackslash warned} \mono{\textbackslash infoed} \mono{\textbackslash subtled} \mono{\textbackslash shaded}}\\
  These variants temporarily change the accent colour.

  \medskip
  \infoed{\structure{Example}}
  \twocolumns{0.5}{
    \begin{block}{Default \mono{block} environment}
      \begin{itemize}
      \item Uses the \shade{\mono{structure}} colour
      \end{itemize}
    \end{block}
  }{%
    \shaded{
      \begin{block}{The exact same \mono{block}}
        \begin{itemize}
        \item But inside \shade{\mono{\textbackslash shaded}}!
        \end{itemize}
      \end{block}
    }
  }

  \begin{lstlisting}[
      basicstyle=\ttfamily\footnotesize,
      language={[LaTeX]TeX},
    ]
\shaded{
  \begin{block}{The exact same \mono{block}}
    \begin{itemize}
    \item But inside \shade{\mono{\textbackslash shaded}}!
    \end{itemize}
  \end{block}
}
  \end{lstlisting}

  \setfootline{%
    These variants essentially replace \mono{exampleblock}/\mono{alertblock}\\
    Note: all elements are affected (e.g., \mono{itemize}'s bullets)
  }
\end{frame}

\begin{frame}[fragile]{Slide Structure}
  \twocolumns{0.42}{
    \structure{Header}\\
    With \mono{\textbackslash structure}
  }{
    \begin{lstlisting}[
        basicstyle=\ttfamily\footnotesize,
        language={[LaTeX]TeX},
      ]^^J
      \\structure\{Header\}\\\\^^J
      With \\mono\{\\textbackslash structure\}^^J
    \end{lstlisting}
  }

  \medskip
  \twocolumns{0.42}{
    \begin{structured}{Header}{Follow-up}
      With \mono{\textbackslash structured} \warn{(new)}
    \end{structured}
  }{
    \begin{lstlisting}[
        basicstyle=\ttfamily\footnotesize,
        language={[LaTeX]TeX},
      ]^^J
      \\begin\{structured\}\{Header\}\{Follow-up\}^^J
      \ \ With \\mono\{\\textbackslash structured\}^^J
      \ \ \\warn\{(new)\}^^J
      \\end\{structured\}^^J
    \end{lstlisting}
  }

  \medskip
  \twocolumns{0.42}{
    \begin{block}{Header}
      With \mono{block}
    \end{block}
  }{
    \begin{lstlisting}[
        basicstyle=\ttfamily\footnotesize,
        language={[LaTeX]TeX},
      ]^^J
      \\begin\{block\}\{Header\}^^J
      \ \ With \\mono\{block\}^^J
      \\end\{block\}^^J
    \end{lstlisting}
  }
\end{frame}


\begin{frame}[fragile]{Automated Side-by-Side Columns \warn{(new)}}
  \structure{Two columns}\\
  Use \shade{\mono{\textbackslash twocolumns}} to create two automatically spaced columns.

  \smallskip
  \twocolumns{0.5}{
    \fbox{\makebox[\textwidth]{ratio=0.5 (provided)}}
  }{
    \fbox{\makebox[\textwidth]{ratio=0.5 (auto)}}
  }
  \smallskip
  \twocolumns{0.66}{
    \fbox{\makebox[\textwidth]{ratio=0.66 (provided)}}
  }{
    \fbox{\makebox[\textwidth]{ratio=0.34 (auto)}}
  }

  \medskip
  \infoed{\structure{Example}}
  \twocolumns{0.66}{
    \begin{block}{Header} Left \end{block}
  }{
    \begin{block}{Header} Right \end{block}
  }
  \begin{lstlisting}[
      basicstyle=\ttfamily\footnotesize,
      language={[LaTeX]TeX},
    ]
\twocolumns{0.66}{
  \begin{block}{Header} Left \end{block}
}{
  \begin{block}{Header} Right \end{block}
}
  \end{lstlisting}
\end{frame}

\begin{frame}[fragile]{Automated Side-by-Side Columns \warn{(new)}}
  \structure{Three columns}\\
  Use \shade{\mono{\textbackslash threecolumns}} to create three automatically spaced columns.

  \smallskip
  \threecolumns{0.25}{0.5}{
    \fbox{\makebox[\textwidth]{0.25 (provided)}}
  }{
    \fbox{\makebox[\textwidth]{0.5 (provided)}}
  }{
    \fbox{\makebox[\textwidth]{0.25 (auto)}}
  }
  \smallskip
  \threecolumns{0.33}{0.34}{
    \fbox{\makebox[\textwidth]{0.33 (provided)}}
  }{
    \fbox{\makebox[\textwidth]{0.34 (provided)}}
  }{
    \fbox{\makebox[\textwidth]{0.33 (auto)}}
  }

  \medskip
  \infoed{\structure{Example}}
  \threecolumns{0.25}{0.5}{
    \begin{block}{Header} Left \end{block}
  }{
    \begin{block}{Header} Middle \end{block}
  }{
    \begin{block}{Header} Right \end{block}
  }
  \begin{lstlisting}[
      basicstyle=\ttfamily\footnotesize,
      language={[LaTeX]TeX},
    ]
\threecolumns{0.25}{0.5}{
  \begin{block}{Header} Left \end{block}}{
  \begin{block}{Header} Middle \end{block}}{
  \begin{block}{Header} Right \end{block}}
  \end{lstlisting}
\end{frame}

\begin{frame}[fragile]{Automated Side-by-Side Columns \warn{(new)}}
  \structure{Vertical alignment}\\
  Use \shade{\mono{\textbackslash twocolumn[bctT]\{\}\{\}}} as you would with \shade{\mono{\textbackslash begin\{columns\}[bctT]}}.

  \medskip
  \structure{Column spacing}\\
  Redefine \shade{\mono{\textbackslash twocolumnsep}}/\shade{\mono{\textbackslash threecolumnsep}} to change column spacing.

  \medskip
  \structure{Long lines}\\
  \oneline{Use \shade{\mono{\textbackslash oneline}} to disable wrapping and stay on a single line.}

  \medskip
  \infoed{\structure{Example}}
  \setlength{\twocolumnsep}{10em}
  \twocolumns{0.5}{
    \fbox{\makebox[\textwidth]{\oneline{left column (very very long)}}}}{
    \fbox{\makebox[\textwidth]{right column}}}
  \vspace{-0.75em}
  \par
  \begin{lstlisting}[
      basicstyle=\ttfamily\footnotesize,
      language={[LaTeX]TeX},
    ]
\setlength{\twocolumnsep}{10em} % default: 2em
\twocolumns{0.5}{
  \fbox{\makebox[\textwidth]{\oneline{left column (very very long)}}}}{
  \fbox{\makebox[\textwidth]{right column}}}
  \end{lstlisting}
  \setfootline{Note: \shade{\mono{\textbackslash oneline}} is very useful to avoid text in \shade{\mono{itemize}} from wrapping}
\end{frame}

\begin{frame}[fragile]{Footline \warn{(new)}}
  \structure{Static footline}\\
  Use \shade{\mono{\textbackslash setfootline}} to add text at the bottom of the slide.\\
  Footnotes reset every new \shade{\mono{frame}}.

  \medskip
  \structure{Multi-line footline}\\
  Simply use \shade{\mono{\textbackslash\textbackslash}} or \shade{\mono{\textbackslash newline}} for a line break, or separate as paragraphs.\\
  %% Do not use \shade{\mono{\textbackslash linebreak}} to avoid justified text.

  \medskip
  \infoed{\structure{Example}}
  \begin{lstlisting}[
      basicstyle=\ttfamily\footnotesize,
      language={[LaTeX]TeX},
    ]
\setfootline{%
  This is a footline\\
  $^1$With a (manually numbered) second line
}
  \end{lstlisting}

  \setfootline{%
    This is a footline\\
    $^1$With a (manually numbered) second line
  }
\end{frame}

\setleftlogo{\includegraphics[height=4ex]{logo-univ-rennes}}
\setrightlogo{\includegraphics[height=4ex]{logo-irisa}}
\begin{frame}[fragile]{Footline Graphics \warn{(new)}}
  \structure{Two possible logo placements}
  \begin{itemize}
  \item \mono{\textbackslash setleftlogo\{\}}
  \item \mono{\textbackslash setrightlogo\{\}}
  \end{itemize}

  \smallskip
  The footline width is automatically recomputed.

  \medskip
  \infoed{\structure{Example}}
  \begin{lstlisting}[
      basicstyle=\ttfamily\footnotesize,
      language={[LaTeX]TeX},
    ]
\setleftlogo{\includegraphics[height=2em]{logo-univ-rennes}}
\setrightlogo{\includegraphics[height=2em]{logo-irisa}}
  \end{lstlisting}

  \setfootline{%
    \faArrowLeft~\shade{\mono{\textbackslash setleftlogo}}
    \hfill
    \shade{\mono{\textbackslash setrightlogo}}~\faArrowRight\\
    \phantom.\hfill page numbers are moved left as needed~\faArrowRight
  }
\end{frame}

\setleftlogo{}
\setrightlogo{}

\begin{frame}{Overlay Page Numbers \warn{(new)}}
  \structure{Don't use \mono{\textbackslash pause}}
  \begin{itemize}[<+->]
  \item screen refreshes are distracting
  \item they waste a lot of attention
  \item they encourage reading the new content out loud
  \item the very last bit of content generally disappears much too fast
  \end{itemize}
  \pause[\thebeamerpauses]

  \bigskip
  \structure{Animation counter}\\
  For the few good uses of overlays, slide numbers are shown as ``\subtle{X$\boldsymbol{\cdot}$Y}''.
  \setfootline{\hfill notice~\faArrowRight}
\end{frame}

\begin{frame}[fragile]{Curated Colour Scheme \warn{(new)}}
  \newcommand\auxpaintdrop[5]{%
    \fill[{#3},draw=#3!75!black,thick,shift={(2.5*#1,2.1*#2)}]
    (-90:0.55) arc (-90:-300:0.55) [out=-30, in=-80]
    to (42:0.85) [out=-60, in=0] to cycle;
/    \node[{#4},style={font=\footnotesize}] at (2.5*#1,2.1*#2-0.8) {\phantom{(}#3\phantom{)}};
    \node[{#3},style={font=\footnotesize}] at (2.5*#1,2.1*#2-1.15) {#5};
  }
  \newcommand\paintdrop[4]{\auxpaintdrop{#1}{#2}{#3}{#3}{#4}}
  \newcommand\paintdropblack[4]{\auxpaintdrop{#1}{#2}{#3}{black}{#4}}

  Base \mono{xcolor} colours are \textcolor{red}{not} \textcolor{green}{the} \textcolor{blue}{most} \textcolor{yellow}{readable} \textcolor{cyan}{or} \textcolor{magenta}{best} \textcolor{brown}{looking}...

  \medskip
  \centerline{
    \begin{tikzpicture}[]
      \paintdrop{0}{0}{lavenderPurple}{(\mono{\textbackslash subtle})}
      \paintdrop{1}{0}{wisteriaViolet}{}
      \paintdrop{2}{0}{cherryPink}{}
      \paintdropblack{3}{0}{penguinWhite}{}
      \paintdrop{4}{0}{ashGray}{}

      \paintdrop{0}{1}{teaBrown}{}
      \paintdrop{1}{1}{crimsonRed}{(\mono{\textbackslash alert})}
      \paintdrop{2}{1}{vividOrange}{(\mono{\textbackslash warn})}
      \paintdrop{3}{1}{goldenYellow}{}
      \paintdrop{4}{1}{steelGray}{(\mono{\textbackslash shade})}

      \paintdrop{0}{2}{lapisBlue}{(\mono{\textbackslash accent})}
      \paintdrop{1}{2}{skyBlue}{}
      \paintdrop{2}{2}{aquaGreen}{}
      \paintdrop{3}{2}{appleGreen}{(\mono{\textbackslash info})}
      \paintdrop{4}{2}{inkBlack}{}
    \end{tikzpicture}
  }

  \setfootline{Note: are also defined \textcolor{steelGray}{steelGr\alert{e}y} and \textcolor{ashGray}{ashGr\alert{e}y}}
\end{frame}


%================================================================================
\appendix
\section{Appendix}
%================================================================================

\setbeamercolor{structure}{fg=steelGray}

\frame[noframenumbering,plain]{%
  \emph{This slide intentionally left blank}%
}

\begin{frame}[fragile]{Appendix}
  \structure{Separated slide numbers}\\
  Simply use \shade{\mono{\textbackslash appendix}} before appendix slides.\\
  Appendix slide numbers reset and start at ``\subtle{+1}''.

  \medskip
  \structure{Separated accent colour}\\
  Redefine Beamer's \shade{\mono{structure}} colour.

  \medskip
  \infoed{\structure{Example}}
  \begin{lstlisting}[
      basicstyle=\ttfamily\footnotesize,
      language={[LaTeX]TeX},
    ]
\appendix
\setbeamercolor{structure}{fg=steelGray}
  \end{lstlisting}

  \setfootline{\hfill notice~\faArrowRight{}}
\end{frame}

\end{document}
